\documentclass[11pt]{article}
\usepackage[a4paper,margin=2.35cm]{geometry}
\usepackage{fontspec}
\usepackage{amsmath,amssymb}
\setmainfont[ItalicFont=Noto Serif Italic]{Malgun Gothic}
\setsansfont{Malgun Gothic}
\setlength{\parindent}{1.2em}
\setlength{\parskip}{0.45em}
\setlength{\emergencystretch}{3em}
\providecommand{\mS}{\mathfrak{S}}
\providecommand{\mA}{\mathfrak{A}}
\providecommand{\srcspaced}[1]{\textbf{#1}}
\begin{document}

이제 $\mS$가 $P_0$에서 완전가약이라고 가정하자. 또한 $P_0$를 완전체라고 가정하면, $P_0$의 모든 대수적 확대 $P$도 마찬가지로 완전체이다. 표현류들은 $P_0$의 임의의 확대 $P$에서도 계속 완전가약이다.\footnote{$P_0$가 완전체가 아닐 때에는 중심의 성분 $K$들이 $P_0$의 ``제1종 확대''가 되는 경우에, 그리고 오직 그 경우에만 완전가약성이 보존된다. 이 가정 아래에서는 본문의 이후 모든 귀결이 성립한다.}

$\mS$의 중심은 유한 개 체의 직합이 된다. 그러한 체 $K$와 동형인 $P_0$의 모든 확대 $P$는 어떤 (단순) 지표의 체가 된다. $P_0$를 $K$와 동형인 $P$로 확대하면, $\mS$에서 하나의 양쪽 단순환이 분리되어 나오며 이 환에는 그 지표의 절대기약 표현류가 대응한다. 이 환은 베더번 정리에 따라 $P$ 위의 한 행렬환과, 동형을 제외하면 유일하게 정해지는 비가환체 $\mA$의 직접곱이 된다. $\mA$의 중심은 $P$이고 $P$ 위에서 차수 $m^2$이며, 여기서 $m$은 그 지표에 속하는 슈어 지수를 뜻한다.

\srcspaced{$\mS$의 모든 분해체---그 지표에 속하는 표현에 관하여---는 동시에 $P$에서의 표현들에 관하여 $\mA$의 분해체이며, 그 역도 성립한다.} $\mS$의 공액 표현들은 $P_0$ 위의 대응하는 양쪽 단순환에서 출발하여 얻는다. 이에 대응하는 체 $\mA_0$는 $\mA$와 동형이 되며 $K$를 중심으로 가진다. 따라서 분해체 문제는 대응하는 비가환체 $\mA$와 그 분해 문제로 완전히 환원된다.\footnote{이는 $P$에서 $\mS$의 기약표현과 가환하는, $P$의 원소를 성분으로 하는 행렬계로 환원하는 Schur의 방법에 대응한다. 이 행렬들의 전치행렬들은 $P$에서 $\mA$의 한 표현을 이룬다.} 특히 다음 정리가 성립한다.

\end{document}
